\chapter*{前言：为什么要把课本“点着”}
\addcontentsline{toc}{chapter}{前言：为什么要把课本“点着”}
\markboth{前言}{}

先说清楚：书名里的“点燃”不是让你真的点火——虽然纸烧掉之后留下灰、放出烟，那本身是另一章要讲的化学反应。我们想点燃的，是你对这个世界的好奇心；想烧掉的，是一种学化学和生物的方式——背。

背元素符号，背方程式，背“线粒体是细胞的能量工厂”，背“显性基因控制显性性状”。背完一场考试，再背下一场。如果你曾经盯着周期表发愣，心里嘀咕“这些东西跟我有什么关系”，这本书就是为你写的。

其实，你每天吃下的面包、手上划破后结的痂、切开后慢慢变褐的苹果、让汽水冒泡的气体，全都是同一套规律在不同尺度上的演出。这套规律可以排成一条长长的因果链：

\begin{quote}
基本粒子 \(\rightarrow\) 原子 \(\rightarrow\) 周期表 \(\rightarrow\) 化学键 \(\rightarrow\) 分子结构 \(\rightarrow\) 化学反应 \(\rightarrow\) 有机分子 \(\rightarrow\) 生物大分子 \(\rightarrow\) 细胞 \(\rightarrow\) 代谢 \(\rightarrow\) DNA \(\rightarrow\) 基因表达 \(\rightarrow\) 遗传 \(\rightarrow\) 进化。
\end{quote}

学校把这条链锯成两段，前半段叫“化学”，后半段叫“生物”，再切成一章一章的知识点。这本书的做法相反：我们沿着链条从头到尾走一遍。你会发现，生命并不是什么“超越化学的神秘力量”——它是一群普普通通的分子，组织到了令人难以置信的程度，以至于能保存过去的信息、维持自身的秩序，还能创造出仍会改变的未来。

走完这一趟，你应该能做到几件以前做不到的事：用电子结构解释周期表为什么长这样，而不是硬背族和周期；用能量和概率判断一个反应“能不能发生”和“有多快”；看懂阿司匹林、咖啡因这些分子的骨架在说什么；分清 DNA、基因、染色体和基因组；面对“保健品抗癌”“基因检测算命”“转基因恐怖故事”这类说法时，有自己的判断力。

我们不敢保证每一页都轻松。有些地方的推理需要你停下来，拿笔画一画、算一算。但我们可以保证三件事：第一，不跳步——每个结论都从你已经理解的东西里长出来；第二，不说谎——暂时还没有答案的问题（比如生命究竟怎样起源），我们会坦白地说“不知道”；第三，不让你白背——每一个符号登场之前，你都会先看到它背后那个真实的现象。

好奇心一旦被点着，课本里那些干巴巴的句子，在你眼里就会全都变成活生生的故事。现在，点火吧。

\hfill 作者
